%% Sample LaTeX file for ePiX                September, 2004
%%
%% Andrew D. Hwang  <ahwang@mathcs.holycross.edu>
%% Department of Mathematics and Computer Science
%% College of the Holy Cross
%% Worcester, MA 01610-2395, USA
%%
\documentclass[11pt]{article}
\usepackage[leqno]{amsmath}
\usepackage{latexsym,pstcol,epic,eepic,rotating}

\newcommand{\ePiX}{\texttt{ePiX}}
\newcommand{\C}{\texttt{C}}
\newcommand{\R}{\mathbf{R}}

\DeclareMathOperator{\re}{Re}
\DeclareMathOperator{\im}{Im}

\setlength{\textheight}{9.5in}
\setlength{\textwidth}{6in}
\setlength{\oddsidemargin}{0.25in}
\setlength{\topmargin}{0in}

\title{\ePiX\ Sample Document}
\author{Version 1.0}
\date{September, 2004}

\begin{document}

\maketitle

\section{Overview}

\ePiX\ is a powerful, flexible, lightweight utility for creating
mathematically accurate \LaTeX\ plots and figures from simple,
mnemonic commands. A detailed user's manual is available in several
formats from
\begin{verbatim}
http://math.holycross.edu/~ahwang/current/ePiX.html
\end{verbatim}
This sample document demonstrates a few of \ePiX's capabilities, with
side-by-side comparisons of input files\footnote{The sample files
distributed with \ePiX\ have \texttt{offset} lines to place the figures
next to their source code.} and corresponding output. \ePiX\ has several 
distinguishing features:
\begin{itemize}

\item Scenes are described in mathematically natural Cartesian
coordinates, making \ePiX\ effectively a vector format. Well-designed
figures are of camera quality over a range of sizes and aspect ratios.
There are almost no default choices; aspect ratio, color, viewpoint
(for 3-D figures) and line width are wholly controllable.

\item High-quality typography is easily incorporated.

\item All the power of \texttt{C++} is available in using and extending 
the program's capabilities.

\item The license, the GNU GPL, is similar to the terms on theorems:
You may run \ePiX\ for any purpose, examine the code, study how the
program works, make improvements, and distribute your improvements so
long as you do not restrict the rights of others to do the same.

\end{itemize}

\section{Examples}

\ePiX\ provides traditional plotting capabilities, such as Cartesian
and polar plots, data plotting from a file, and parametric curves and
surfaces (without automatic hidden object removal). As mentioned, the 
implementation allows \ePiX\ to be regarded as a programming language;
any numerical algorithm written in \texttt{C}~or \texttt{C++} can be 
used in an \ePiX\ figure. The use of programming constructs can make a
figure flexible (so that the appearance can be precisely but dramatically
altered by making a few small changes to the input file) and easier to
maintain.

The figures below were chosen to emphasize results that are relatively
difficult to achieve with existing Free (and some commercial) plotting
software.  Simple line figures containing polygons, ellipses, and
splines are almost self-explanatory, so no examples are given. 

In the files below, several short statements are often put on a single
line to make the file fit on a page; this is not good programming
practice, and should be avoided in real files. Several sample files
could be shortened by hard-wiring constants.

\clearpage

\begin{figure}[hbt]
\begin{center}
\input{parabola.eepic}
\vspace*{-2in}
\begin{verbatim}
#include "epix.h"
using namespace ePiX;

// the function to be plotted
double f(double x) { return x*x; }

int main() 
{
  bounding_box(P(-2,0),P(2,4));
  unitlength("1in");
  picture(3,3);

  begin();

  grid(32,32); // fine grid

  bold();
  grid(4,4); // and coarse

  h_axis_labels(4, P(0,-2), b);
  v_axis_masklabels(P(0, 1), P(0, y_max), 3, P(0,0), c); // omit 0

  red();
  plot(f,x_min,x_max,20);
  
  end();
}
\end{verbatim}
\caption{A basic plot.}
\label{fig:parabola}
\end{center}
\end{figure}

\clearpage

\begin{figure}[hbt]
\begin{center}
\input{cropplot.eepic}
\vspace*{-200pt}
\begin{small}
\begin{verbatim}
#include "epix.h"
using namespace ePiX;

double f(double t) { return 2*t*(1-t)*(1-t); }
double g(double t) { return 1/(1-t*t); }

int main() 
{
  bounding_box(P(-2,-4), P(2,4));
  picture(200,200);
  unitlength("1pt");

  begin();

  crop(); 
  // Vertical asymptotes
  dashed();
  line(P(-1, y_min), P(-1, y_max));
  line(P( 1, y_min), P( 1, y_max));
  solid();

  // Axes
  h_axis(8);
  v_axis(8);

  h_axis_labels(4, P(-1, 2), tl); // align top-left
  v_axis_labels(4, P(-1, 2), tl);

  // Graphs
  plot(f, x_min, x_max, 40);
  bold();
  plot(g, x_min, x_max, 80);

  end();
}

\end{verbatim}
\end{small}
\caption{``Cropping'' removes elements outside the bounding box.}
\label{fig:cropplot}
\end{center}
\end{figure}

\clearpage

\noindent There are angular ``modes'' for polar and spherical
coordinates. \ePiX\ defines its own trig functions, which are
sensitive to the current mode.

\begin{figure}[hbt]
\begin{center}
\input{polar.eepic}
\vspace*{-150pt}
\begin{verbatim}
#include "epix.h"
using namespace ePiX;

double f(double t) 
{ 
  return 2*Cos(3*t); 
}

int main() 
{
  unitlength("1pt");
  bounding_box(P(-2,-2), P(2,2));
  picture(180, 180);

  begin();
  degrees();

  h_axis(P(x_min, y_min), P(x_max, y_min), 8);
  v_axis(P(x_min, y_min), P(x_min, y_max), 8);
  
  polar_grid(2, 4, 24); // radius, rings, sectors

  h_axis_labels(P(x_min,y_min), P(x_max,y_min), x_size, P(-12,-14));
  v_axis_labels(P(x_min,y_min), P(x_min,y_max), y_size, P(-4,0), l);

  bold();
  polarplot(f, 0, 180, 120); // plot over [0, 180] using 120 intervals

  end();
}

\end{verbatim}
\caption{The polar graph $r=2\cos 3\theta$ for $0\leq\theta\leq\pi$.}
\label{fig:polar}
\end{center}
\end{figure}

\clearpage

\noindent Regions between graphs can be shaded; gray density ranges
from~0 (white) to~1 (black).

\begin{figure}[hbt]
\begin{center}
\input{shadeplot.eepic}
\vspace*{-170pt}
\begin{footnotesize}
\begin{verbatim}
#include "epix.h"
using namespace ePiX;
double k=4;           // change width of hump
double dx = 0.05;     // width of thin shaded region
double x = 1/sqrt(k); // position of thin shaded region
double dy = 0.25;     // for vertical distance between labels; kludgey

double f(double t) { return sqrt(fabs(k)/(2*M_PI))*exp(-k*t*t); }
P pt1 = P(x, f(x)+2*dy), pt2 = P(x+dx,f(x)+dy), pt3 = P(x+3*dx,f(x));

int main() 
{
  unitlength("1pt");
  picture(150, 150);
  bounding_box(P(0,0), P(1,1));

  begin();
  fill(); 
  gray(0.1);     shadeplot(f, x_min, x, 90);
  gray(0.4);     rect(P(x,f(x)), P(x+dx, 0));
  gray(0.6);     shadeplot(f, x, x+dx, 10);
  fill(false);

  arrow_camber(0.25);  arrow_fill(1);
  arrow(pt1, P(0.5*(x_min+x), 0.5*f(0.5*(x_min+x))),
        P(0,0), "$F(x)=\\displaystyle\\int_a^x f(t)\\,dt$", tr);

  arrow(pt2, P(x+0.5*dx, f(x)), 
        P(0,0), "Area of rectangle = $f(x)\\,dx$", tr);

  arrow(pt3, P(x+dx, 0.5*f(x+dx)), 
        P(0,0), "Area = $F(x+dx)-F(x)$", tr);

  bold();
  plot(f, x_min, x_max, 120);

  plain(); h_axis(4); v_axis(4);

  label(P(x_min,0), P(0,-5), "$a$", b);
  label(P(x,0), P(0,-5), "$x$", b);
  label(P(x+dx,0), P(0,-2), "$x+dx$", br);

  end();

}
\end{verbatim}
\end{footnotesize}
\caption{The first fundamental theorem of calculus: $F'(x)=f(x)+o(1)$.}
\label{fig:shadeplot}
\end{center}
\end{figure}

\clearpage

\noindent Control structures---loops and decision statements---can be
used to create input files whose logical structure matches the
mathematical structure of the figure.

\begin{figure}[hbt]
\begin{center}
\input{newton.eepic}
\vspace*{-2.5in}
\begin{footnotesize}
\begin{verbatim}
#include "epix.h"
using namespace ePiX;

double f(double t) { return t*t*t-3*t+1; }
double df(double t) { return 3*t*t - 3; }

int main() 
{
  bounding_box(P(1.5,0),P(2,3));
  unitlength("1in");
  picture(2.5,2.5);

  begin();
  h_axis(2);

  double x0 = 2;
  label(P(x0,0), P(0,-4), "$x_i$", b);

  for (int i=0; i < 3; ++i)
    {
      dashed();     line(P(x0,0), P(x0,f(x0)));
      solid();      line(P(x0,f(x0)), P(x0-f(x0)/df(x0),0));

      x0 -= f(x0)/df(x0); // Newton's method
    }

  bold();    plot(f, x_min, x_max+0.05, 60);

  double x1 = 2-f(2)/df(2);

  label(P(1.75, f(1.75)), P(-2,2), "$y=f(x)$", tl);
  label(P(x1,0), P(0,-4), "$x_{i+1}$", b);
  label(P(1.75, df(2)*(1.75-x1)), P(0,-4), "Slope $= f'(x_i)$", br);

  end();
}
\end{verbatim}
\end{footnotesize}
\caption{Newton's method for root approximation}
\label{fig:newton}
\end{center}
\end{figure}

\clearpage

\noindent\ePiX\ can graph derivatives and definite integrals:

\begin{figure}[hbt]
\begin{center}
\input{calculus.eepic}
\vspace*{-120pt}
\begin{small}
\begin{verbatim}
#include "epix.h"
using namespace ePiX;

double MAX=2*M_PI;
double f(double t)
{
  return t*Sin(t);
}

int main() 
{
  unitlength("1pt");
  picture(240, 120);
  bounding_box(P(-MAX,-MAX), P(MAX,MAX));

  begin();

  // Coordinate axes and labels
  h_axis(8);
  v_axis(4);
  font_size("scriptsize");

  label(P(0,y_max), P(-4,0), "$2\\pi$", l);
  label(P(0,y_min), P(-4,0), "$-2\\pi$", l);

  label(P(x_min,0), P(0,2), "$-2\\pi$", t);
  label(P(x_max,0), P(0,2), "$2\\pi$", t);

  bold();
  plot(f, x_min, x_max, 90);
  green();
  plot(D(f), x_min, x_max, 90);    // derivative class "D"

  blue();
  plot(I(f, 0), x_min, x_max, 90); // definite integral from x=0

  end();
}
\end{verbatim}
\end{small}
\caption{$y=x\sin x$ (black), its derivative (green) and integral from~$0$
(blue).}
\label{fig:calculus}
\end{center}
\end{figure}

\clearpage

\noindent Solutions of ODEs are computed with Euler's method. Vector 
fields may be plotted at constant (shown) or true length.

\begin{figure}[ht]
\begin{center}
\input{slopefield.eepic}
\vspace*{-180pt}
%%\begin{footnotesize}
\begin{verbatim}
#include "epix.h"
using namespace ePiX;

P F(double s, double t)
{
  return (P(0.1,0.025)&P(s,t)) + 
    (1/(0.01+s*s+t*t))*P(-t,s);
}

int main() 
{
  unitlength("1pt");
  bounding_box(P(-4, -3), P(2,2));
  picture(234, 195);

  begin();
  crop();

  blue(0.4);
  dart_field(F, P(x_min, y_min), P(x_max, y_max), 4*x_size, 4*y_size);

  bold();
  for (int i=0; i<7; ++i)
    {
      rgb(0.25 - 0.05*i, 1 - 0.1*i, 0.2*i);
      ode_plot(F, P(-0.9-0.025*i,0), 20, 200);
    }
  end();
}

\end{verbatim}
%%\end{footnotesize}
\caption{A slope field and six solutions of an ODE.}
\label{fig:slopefield}
\end{center}
\end{figure}

\clearpage

\noindent The ``denominator'' function~$f$, defined by
\begin{equation*}
f(x)=\begin{cases}
\frac{1}{q} & \text{if $x=\frac{p}{q}$ in lowest terms} \\
0 & \text{if $x$~is irrational}
\end{cases}
\end{equation*}
can be plotted with a nested \texttt{for} loop.

\begin{figure}[hbt]
\begin{center}
\input{denom.eepic}
\vfil

\begin{verbatim}
#include "epix.h"
using namespace ePiX;

int N = 30; // maximum denominator plotted

int main() 
{
  unitlength("0.01in");
  bounding_box(P(-2,0), P(2,1));
  picture(400,100);

  begin();   

  h_axis(2*x_size);
  v_axis(2);

  h_axis_labels(x_size, P(-12, -12));

  dot_size(1); // draw dots 1pt in diameter
  for (int i=1; i< N; ++i) 
    for (int j=i*x_min; j <= i*x_max; ++j) 
      if (gcd(i, j) == 1) 
        dot(P(j*1.0/i, 1.0/i));

  end();
}

\end{verbatim}
\caption{A "pathological" function in real analysis.}
\label{fig:denom}
\end{center}
\end{figure}

\clearpage

\noindent Finite sums are defined by an algorithm and are therefore
easy to plot. The function being scaled and summed is~cb, the
``Charlie Brown'' function (blue). Nowhere differentiability is
essentially obvious from the picture, because the graph is
self-similar (red) and not a line.

\begin{figure}[hbt]
\begin{center}
\input{weierstrass.eepic}
\vspace*{-1.375in}
\begin{small}
\begin{verbatim}
#include "epix.h"
using namespace ePiX;

const int N=8; // Number of summands

double weierstrass(double t)
{
  double y=0;
  for(int i=0; i < N; ++i)
    y += pow(2,-i)*cb(pow(2,i)*t);

  return y;
}

int main() 
{
  bounding_box(P(-2, 0), P(2, 1.5));
  picture(320, 120);
  unitlength("0.01in");

  begin();

  h_axis(2*x_size);
  v_axis(2*y_size);
  h_axis_labels(x_size, P(-12,-14));
  
  blue();
  plot(cb, x_min - 0.25, x_max+0.25, 4*x_size + 2);

  bold();
  black();
  plot(weierstrass, x_min, x_max, pow(2,N));

  red();
  plot(weierstrass, 0.5, 1.5, pow(2,N-2));

  end();
}

\end{verbatim}
\end{small}
\caption{A Weierstrass nowhere-differentiable function.}
\label{fig:weierstrass}
\end{center}
\end{figure}

\clearpage

\ePiX\ can approximate the extreme values of a function on an
interval, which is useful for drawing inscribed or circumscribed
rectangles in a graph. The sine function is symbolized by~\texttt{f}
throughout the body (except the label, of course), so re-drawing the
figure with a different integrand only requires changing the
definition of~\texttt{f} and re-sizing the \verb+bounding_box+.  The
number of rectangles is similarly ``parametrized''.

\begin{figure}[hbt]
\begin{center}
\input{uppersum.eepic}
\vspace*{-1in}
\begin{small}
\begin{verbatim}
#include "epix.h"
using namespace ePiX;

const int N=12; // Number of rectangles
double f(double t) { return Sin(t); } // gather references to integrand

int main() 
{
  bounding_box(P(0,0), P(3,1));
  unitlength("1in");
  picture(3, 1);

  begin();

  double dx = x_size/N;
  gray(0.1); // (un)set filling in loop

  riemann_sum(f, x_min, x_max, N, UPPER);
  fill();
  riemann_sum(f, x_min, x_max, N, LOWER);
  fill(false);

  h_axis(x_size);
  v_axis(2*y_size);

  h_axis_labels(x_size, P(0,-4), b);
  label(P(x_max, f(x_max)), P(4,0), "$y=\\sin x$", r);

  bold();    blue();
  plot(f, x_min, x_max, 40);

  end();
}

\end{verbatim}
\end{small}
\caption{Upper and lower sums for $\displaystyle\int_0^3\sin x\,dx$.}
\label{fig:uppersum}
\end{center}
\end{figure}

\clearpage

A loop index may be used to control an entire figure, generating a
sequence of snapshots of a time-varying picture, such as a rolling
wheel, or the flow of an ODE. The \ePiX\ project page contains links
to animations that can be played in any web browser.

\begin{small}
\begin{verbatim}
#include "epix.h"
using namespace ePiX;

P F(double t, double r)
{
  return P(t-r*Sin(t), 1-r*Cos(t));
}

int main() 
{
  const double dt = 5*M_PI/11;
  double t = 0;

  picture(8, 1);
  bounding_box(P(-1,0), P(15,2));
  unitlength("0.625in");

  for(int i=0; i < 9; ++i)
    {
      t += dt;
      begin(); // Entire picture inside loop body

      line(P(x_min, y_min), P(x_max, y_min)); // the ground

      circle wheel(P(t,1), 1);  wheel.draw();

      domain R = domain(P(0,0), P(t, 1), // [0,t] x [0,1]
                        mesh(10*i,5), mesh((int) ceil(1+4*t), 5));

      // the paths
      bold();
      for (int j=0; j < 6; ++j)
        {
          rgb(1-0.125*j, 0.125*j, 0.5+0.25*j);
          plot(F, R.slice2(0.2*j));
        }

      // the spoke
      green();
      line(P(t,1), F(t,1));

      end();
    
      // figure separator and vertical space
      if (i < 8)
        printf("\n\\vspace*{3ex}\n%%%%");
    }
}

\end{verbatim}
\end{small}

\vspace*{-100pt}
\begin{figure}[hb]
\begin{center}
\input{wheel.eepic}
\caption{Snapshots of cycloids.}
\label{fig:wheel}
\end{center}
\end{figure}

\clearpage

Figure~\ref{fig:sphere} (by Jacques L'helgoual) demonstrates some
spherical geometry capabilities: Space curves can be projected to the
unit sphere radially, and plane curves can be projected
stereographically from the north or south pole, with or without hidden
line removal.

\begin{figure}[hbt]
\begin{center}
\input{sphere.eepic}
\vspace*{-150pt}
\begin{scriptsize}
\begin{verbatim}
#include "epix.h"
using namespace ePiX;

const double k = 2*M_PI/(360*sqrt(3)); // assume "degrees" mode
double exp_cos(double t) { return exp(k*t)*Cos(t); }
double exp_sin(double t) { return exp(k*t)*Sin(t); }
double minus_exp_cos(double t) { return -exp_cos(t); }
double minus_exp_sin(double t) { return -exp_sin(t); }

int main()
{
  bounding_box(P(-1,-1), P(1,1));
  picture(160,160);
  unitlength("1pt");

  begin();
  degrees(); // set angle units
  viewpoint(1, 2.5, 3);

  sphere S1=sphere(); // unit sphere
  S1.draw();

  pen(0.15); // hidden portions of loxodromes
  blue();
  backplot_N(exp_cos, exp_sin, -540, 540, 90);
  backplot_N(minus_exp_cos, minus_exp_sin, -540, 540, 90);
  red();
  backplot_N(exp_sin, minus_exp_cos, -540, 540, 90);
  backplot_N(minus_exp_sin, exp_cos, -540, 540, 90);

  black(); // coordinate grid
  for (int i=0; i<=12; ++i) { latitude(90-15*i,0,360); longitude(30*i,0,360); }

  bold(); // visible portions of loxodromes
  blue();
  frontplot_N(exp_cos, exp_sin, -540, 540, 360);
  frontplot_N(minus_exp_cos, minus_exp_sin, -540, 540, 360);
  red();
  frontplot_N(exp_sin, minus_exp_cos, -540, 540, 360);
  frontplot_N(minus_exp_sin, exp_cos, -540, 540, 360);

  end();
}
\end{verbatim}
\end{scriptsize}
\caption{Loxodromes on the unit sphere.}
\label{fig:sphere}
\end{center}
\end{figure}

\clearpage

\noindent \texttt{PSTricks} can be incorporated in \ePiX\ files to 
achieve colored filling and other effects. The \texttt{psset} command
puts its argument into the output file.

\begin{figure}[hbt]
\begin{center}
\input{geomsum.eepic}
\vspace*{-192pt}
%%\begin{small}
\begin{verbatim}
#include "epix.h"
using namespace ePiX;

const int N=8;

int main() 
{
  unitlength("1.5pt");
  bounding_box(P(0,0), P(1,1));
  picture(128,128);

  use_pstricks();
  begin();

  psset("fillcolor=lightgray, linecolor=white");
  rect(P(0,0), P(1,1));
  label(P(1.0/4,   1.0/2), "$\\frac{1}{2}$");
  label(P(5.0/8,   3.0/4), "$\\frac{1}{8}$");
  label(P(13.0/16, 7.0/8), "$\\frac{1}{32}$");

  psset("fillcolor=blue");
  double t=0.5;

  for(int i=0; i<N; ++i, t *= 0.5)
    {
      rect(P(1-t, 1-2*t), P(1, 1-t));
      line(P(1-t, 1-2*t), P(1-t, 1));
    }
  white();
  label(P(3.0/4,     1.0/4), "$\\mathbf{\\frac{1}{4}}$");
  label(P(7.0/8,     5.0/8), "$\\mathbf{\\frac{1}{16}}$");
  label(P(15.0/16, 13.0/16), "$\\mathbf{\\frac{1}{64}}$");

  end();
}
\end{verbatim}
%\end{small}
\caption{A geometric series with ratio $1/2$.}
\label{fig:geomsum}
\end{center}
\end{figure}

\clearpage

\noindent \ePiX\ computes intersections of geometric objects such as lines,
circles, and planes.
%% In this file, newlines have been omitted to make the file fit on one 
%% page; generally it is a good idea to have only one statement per line.

\begin{figure}[hbt]
\begin{center}
\input{pascal.eepic}
\vspace*{-2.5in}
\begin{footnotesize}
\begin{verbatim}
#include "epix.h"
using namespace ePiX;

int main() 
{
  bounding_box(P(-2,-2),P(2,2));
  unitlength("1in");
  picture(3,3);

  begin(); 
  crop();


  P V = P(2,-0.25),   W = P(3,1);
  P P1=P(-1.5,-1),    P2 = P1 + V,       P3 = P1 + 1.5*V;    // points
  P Q1 = P(-1,1),     Q2 = Q1 + 0.5*W,   Q3 = Q1 + W;

  segment L12(P1, Q2),   L13(P1, Q3);
  segment L21(P2, Q1),   L23(P2, Q3);
  segment L31(P3, Q1),   L32(P3, Q2);

  P R1 = L12*L21,    R2 = L13*L31,    R3 = L32*L23;   // points of intersection

  dot(P1, P(0,-2), "$P_1$", b);   dot(Q1, P(0,2), "$Q_1$", t); 
  dot(P2, P(0,-2), "$P_2$", b);   dot(Q2, P(0,2), "$Q_2$", t);
  dot(P3, P(0,-2), "$P_3$", b);   dot(Q3, P(0,2), "$Q_3$", t);   

  red();  L12.draw();  L21.draw();  L13.draw();  L31.draw();  L32.draw();  L23.draw();

  green();    Line(P1,P3);    Line(Q1,Q3);

  blue();
  box(R1, P(4,2), "$p_1$", r);
  box(R2, P(4,2), "$p_2$", r);
  box(R3, P(4,2), "$p_3$", r);

  dashed();    Line(R1,R3);
  end();
}
\end{verbatim}
\end{footnotesize}
\caption{Pascal's theorem on sides of a hexagon.}
\label{fig:pascal}
\end{center}
\end{figure}

\clearpage

\noindent Paths may be built in pieces and manipulated as a single object.

\begin{figure}[hbt]
\begin{center}
\input{contour.eepic}
\vspace*{-2.25in}
\begin{footnotesize}
\begin{verbatim}
#include "epix.h"
using namespace ePiX;

int main() 
{
  bounding_box(P(-5, -5), P(5,5));
  picture(160, 160);
  unitlength("0.35mm");

  use_pstricks();
  begin();  degrees();
  double theta=60;
  double rad=0.25, Rad=4.5; // arc radii
  double ht=1.0/16;         // half-width of slit

  // "start" angles of arcs
  double theta1 = Asin(ht/rad), theta2 = Asin(ht/Rad);

  // build keyhole in pieces; += concatenates, -= reverses orientation
  path contour(P(0,0), Rad*E_1, Rad*E_2, theta2, 360-theta2);    // outer arc
  contour -= path(polar(rad,-theta1), polar(Rad,-theta2));       // lower slit
  contour -= path(P(0,0), rad*E_1, rad*E_2, theta1, 360-theta1); // inner arc
  contour += path(polar(rad, theta1), polar(Rad, theta2));       // upper slit

  fill();
  std::cout << "\n\\newrgbcolor{orange}{1 0.7 0.2}";

  psset("fillcolor=orange,linecolor=red,linewidth=1.5pt");
  contour.set_fill();
  contour.draw();
  use_pstricks(false);

  dot(P(0,0));   // the origin
  gray(1); // blacken arrowheads
  arrow(P(Rad/4, 0.1*Rad), P(3*Rad/4, 0.1*Rad)); // arrow ends in terms of Rad
  arrow(P(3*Rad/4, -0.1*Rad), P(Rad/4, -0.1*Rad));
  arc_arrow(P(0,0), 0.9*Rad, 180-theta, 180+theta);

  label(P(Rad,ht), P(2,4), "$R\\to\\infty$", tr);
  label(P(0,rad), P(0,4), "$\\delta\\to0$", tl);
  label(polar(Rad, 45), P(0,0), "$\\gamma$", tr);

  end();
}
\end{verbatim}
\end{footnotesize}
\caption{A keyhole contour for a branch cut integral.}
\label{fig:contour}
\end{center}
\end{figure}

\clearpage

\noindent \ePiX\ provides ``clipping'', which removes figure elements
that lie outside a specified box and can be used (with manual
ordering) to layer figure elements for hidden object removal.

\begin{figure}[hbt]
\begin{center}
\input{pole.eepic}
\vspace*{-2.5in}
\begin{verbatim}
#include "epix.h"
using namespace ePiX;
const int M=8;
double MAX=1.5, Bd=2;
P F(double x, double y) // fake pole
{ return P(x, y, x/(0.01+x*x+y*y)); }
P pt1 = P(-MAX,-MAX), pt2 = P(MAX, MAX);

int main() 
{
  bounding_box(P(-Bd,-Bd), P(Bd,Bd));
  unitlength("1in");
  picture(2.5,2.5);

  begin();
  domain R(pt1, pt2, mesh(2*M,2*M), mesh(4*M,4*M));
  label(P(x_min, y_max), P(2,-2), "$z=\\mathrm{Re}\\,\\frac{1}{x+iy}$", br);
  viewpoint(6,8,5); Light=P(-4,10,10);

  clip(); clip_box(P(-Bd,-Bd,-Bd), P(Bd,Bd,0)); fill(); surface(F, R); // bottom half
  gray(0); rect(pt1, pt2);  fill(false);
  pen(0.15); rgb(0.75,0.75,0.75); plot(F, R); plain();  // simulate transparency

  clip(false); magenta(); grid(pt1, pt2, M, M); // axes
  arrow(P(-Bd,0,0), P(Bd,0,0));  label(P(Bd,0), P(-4,-2), "$x$", l); 
  arrow(P(0,-Bd,0), P(0,Bd,0));  label(P(0,Bd), P( 2,-2), "$y$", br);
  arrow(P(0, 0, 0), P(0,0,2.5)); label(P(0,0,2.5), P(0,4), "$z$", t);

  clip(); clip_box(P(-Bd,-Bd,0), P(Bd,Bd,Bd)); fill(); surface(F, R); // top half
  fill(false); pen(1.5); red(); plot(F, R.resize1(0.5,MAX).slice2(0));
  end();
}
\end{verbatim}
\caption{The real part of the complex reciprocal.}
\label{fig:pole}
\end{center}
\end{figure}

\clearpage

\begin{figure}[hbt]
\begin{center}
\input{sqrt.eepic}
\vspace*{-2.25in}
\begin{footnotesize}
\begin{verbatim}
#include "epix.h"
using namespace ePiX;

P f(double x, double y)
{
  pair z = pair(x,y);
  pair temp = z*z;
  return P(temp.x1(), temp.x2(), x);
}

inline P F(double r, double t)
{ return f(r*Cos(t), r*Sin(t)); }

P g(double t) { return t*P(t,0,1); }
const double MAX=1.5;

domain R(P(0,0), P(1.25,0.5), mesh(6,24), mesh(12,60));

int main() 
{
  bounding_box(P(-MAX,-MAX),P(MAX,MAX));
  unitlength("1in");
  picture(2.5,2.5);

  begin();
  revolutions();    viewpoint(4,-2,3);

  pen(0.15);  rgb(0.7,0.7,1); 
  plot(F, R);

  plain();  blue(); 
  plot(F, R.resize2(0.5,1));

  black();  pen(0.5);
  arrow(P(0,0,0), 2*E_1);
  arrow(P(0,0,0), 2*E_2);
  arrow(P(0,0,0), 1.5*E_3);

  masklabel(2*E_1, P(0,0), "$\\mathrm{Re}\\,z$", r);
  label(2*E_2, P(0,0), "$\\mathrm{Im}\\,z$", r);
  label(1.5*E_3, P(2,0), "$\\mathrm{Re}\\,\\sqrt{z}$", r);

  bold(); red();
  plot(g, -1.25, 1.25, 40);
  end();
}
\end{verbatim}
\end{footnotesize}
\caption{Two sheets of the Riemann surface of~$\sqrt{z}$.}
\label{fig:sqrt}
\end{center}
\end{figure}

\clearpage

\subsection*{Gallery}

\noindent The input files for the figures below are in the
\texttt{samples} directory of the source tree.

\begin{center}
\input{lissajous.eepic}\hspace*{1in}
\input{koch.eepic}
\vspace*{0.5in}

\input{oscillator.eepic}\hspace*{1in}
\input{torus.eepic}
\vfil

\input{levelset.eepic}\hspace*{1in}
\input{clipping.eepic}
\end{center}

\end{document}
